\chapter{Unit 5: Fundamentals of Filters and Operational Amplifier}

\section{Section A: Multiple Choice Question Practice (60 Curated Problems)}

\noindent\textbf{Q1. A low-pass filter allows which type of signals to pass through?} \hfill [Topic: low-pass filter | Easy]
\begin{enumerate}[label=(\alph*),topsep=2pt,itemsep=1pt,partopsep=0pt,parsep=0pt]
    \item {Only DC signals}
    \item {Signals within a narrow band in the middle of the spectrum, while attenuating everything above and below that band}
    \item {High-frequency signals}
    \item {Low-frequency signals}
\end{enumerate}
\vspace{2mm}

\noindent\textbf{Q2. A high-pass filter attenuates signals that are:} \hfill [Topic: high-pass filter | Easy]
\begin{enumerate}[label=(\alph*),topsep=2pt,itemsep=1pt,partopsep=0pt,parsep=0pt]
    \item {Equal to the cutoff frequency}
    \item {Above the cutoff frequency}
    \item {At all frequencies equally}
    \item {Below the cutoff frequency}
\end{enumerate}
\vspace{2mm}

\noindent\textbf{Q3. A band-pass filter is designed to:} \hfill [Topic: band-pass filter | Easy]
\begin{enumerate}[label=(\alph*),topsep=2pt,itemsep=1pt,partopsep=0pt,parsep=0pt]
    \item {Pass a specific range of frequencies and reject others}
    \item {Reject all frequencies}
    \item {Pass only DC voltage}
    \item {Pass only the lowest frequency}
\end{enumerate}
\vspace{2mm}

\noindent\textbf{Q4. An ideal operational amplifier has an input impedance that is:} \hfill [Topic: operational amplifier abstraction- device properties of the operational amplifier | Easy]
\begin{enumerate}[label=(\alph*),topsep=2pt,itemsep=1pt,partopsep=0pt,parsep=0pt]
    \item {Zero}
    \item {Negative}
    \item {Equal to 1 , k O m e g a}
    \item {Infinite}
\end{enumerate}
\vspace{2mm}

\noindent\textbf{Q5. An ideal operational amplifier has an output impedance of:} \hfill [Topic: operational amplifier abstraction- device properties of the operational amplifier | Easy]
\begin{enumerate}[label=(\alph*),topsep=2pt,itemsep=1pt,partopsep=0pt,parsep=0pt]
    \item {Zero}
    \item {10 , k O m e g a}
    \item {Equal to the input impedance}
    \item {Infinite}
\end{enumerate}
\vspace{2mm}

\noindent\textbf{Q6. The open-loop voltage gain of an ideal op-amp is:} \hfill [Topic: properties of operational amplifier | Easy]
\begin{enumerate}[label=(\alph*),topsep=2pt,itemsep=1pt,partopsep=0pt,parsep=0pt]
    \item {Infinite}
    \item {Exactly 100}
    \item {Zero}
    \item {Unity}
\end{enumerate}
\vspace{2mm}

\noindent\textbf{Q7. In the virtual ground concept, the voltage at the inverting input of an ideal op-amp with negative feedback is approximately:} \hfill [Topic: simple op amp circuits -- virtual ground concept | Easy]
\begin{enumerate}[label=(\alph*),topsep=2pt,itemsep=1pt,partopsep=0pt,parsep=0pt]
    \item {Equal to the supply voltage}
    \item {0 , V}
    \item {Equal to the output voltage}
    \item {Infinite}
\end{enumerate}
\vspace{2mm}

\noindent\textbf{Q8. The output of an inverting amplifier is:} \hfill [Topic: inverting and non-inverting op-amp | Easy]
\begin{enumerate}[label=(\alph*),topsep=2pt,itemsep=1pt,partopsep=0pt,parsep=0pt]
    \item {Always zero}
    \item {In phase with the input}
    \item {$180^\circ$ out of phase with the input}
    \item {$90^\circ$ out of phase with the input}
\end{enumerate}
\vspace{2mm}

\noindent\textbf{Q9. The voltage gain of an inverting amplifier with input resistor R 1 and feedback resistor R f is:} \hfill [Topic: inverting and non-inverting op-amp | Easy]
\begin{enumerate}[label=(\alph*),topsep=2pt,itemsep=1pt,partopsep=0pt,parsep=0pt]
    \item {- f r a c R f R 1}
    \item {f r a c R 1 R f}
    \item {R f t im es R 1}
    \item {1 + f r a c R f R 1}
\end{enumerate}
\vspace{2mm}

\noindent\textbf{Q10. The voltage gain of a non-inverting amplifier is given by:} \hfill [Topic: inverting and non-inverting op-amp | Easy]
\begin{enumerate}[label=(\alph*),topsep=2pt,itemsep=1pt,partopsep=0pt,parsep=0pt]
    \item {- f r a c R f R 1}
    \item {1 + f r a c R f R 1}
    \item {0}
    \item {f r a c R 1 R f}
\end{enumerate}
\vspace{2mm}

\noindent\textbf{Q11. An op-amp summing amplifier produces an output proportional to:} \hfill [Topic: op-amp as an adder and subtractor | Easy]
\begin{enumerate}[label=(\alph*),topsep=2pt,itemsep=1pt,partopsep=0pt,parsep=0pt]
    \item {The sum of the input voltages}
    \item {The difference of the input voltages}
    \item {The ratio of the input voltages}
    \item {The product of the input voltages}
\end{enumerate}
\vspace{2mm}

\noindent\textbf{Q12. An op-amp subtractor circuit gives an output proportional to the:} \hfill [Topic: op-amp as an adder and subtractor | Easy]
\begin{enumerate}[label=(\alph*),topsep=2pt,itemsep=1pt,partopsep=0pt,parsep=0pt]
    \item {Square of the input voltage}
    \item {Difference between two input voltages}
    \item {Sum of two input voltages}
    \item {Integral of the input voltage}
\end{enumerate}
\vspace{2mm}

\noindent\textbf{Q13. An op-amp integrator produces an output that is the:} \hfill [Topic: op-amp RC circuits -- op-amp integrator | Easy]
\begin{enumerate}[label=(\alph*),topsep=2pt,itemsep=1pt,partopsep=0pt,parsep=0pt]
    \item {Inverse of the input frequency}
    \item {Sum of two inputs}
    \item {Derivative of the input signal}
    \item {Integral of the input signal}
\end{enumerate}
\vspace{2mm}

\noindent\textbf{Q14. In a basic op-amp integrator circuit, the feedback element is a:} \hfill [Topic: op-amp RC circuits -- op-amp integrator | Easy]
\begin{enumerate}[label=(\alph*),topsep=2pt,itemsep=1pt,partopsep=0pt,parsep=0pt]
    \item {Diode}
    \item {Resistor}
    \item {Capacitor}
    \item {Inductor}
\end{enumerate}
\vspace{2mm}

\noindent\textbf{Q15. An op-amp differentiator gives an output proportional to the:} \hfill [Topic: op-amp differentiator | Easy]
\begin{enumerate}[label=(\alph*),topsep=2pt,itemsep=1pt,partopsep=0pt,parsep=0pt]
    \item {Integral of the input signal}
    \item {Rate of change of the input signal}
    \item {Average of the input signal}
    \item {Peak value of the input signal}
\end{enumerate}
\vspace{2mm}

\noindent\textbf{Q16. In a basic op-amp differentiator circuit, the input element is a:} \hfill [Topic: op-amp differentiator | Easy]
\begin{enumerate}[label=(\alph*),topsep=2pt,itemsep=1pt,partopsep=0pt,parsep=0pt]
    \item {Capacitor}
    \item {Resistor}
    \item {Transistor}
    \item {Inductor}
\end{enumerate}
\vspace{2mm}

\noindent\textbf{Q17. An active high-pass filter differs from a passive one because it includes an:} \hfill [Topic: active high pass filter | Easy]
\begin{enumerate}[label=(\alph*),topsep=2pt,itemsep=1pt,partopsep=0pt,parsep=0pt]
    \item {Extra ground connection}
    \item {Operational amplifier for gain}
    \item {Additional transformer}
    \item {External battery only}
\end{enumerate}
\vspace{2mm}

\noindent\textbf{Q18. An active band-stop filter is designed to:} \hfill [Topic: active band stop filter | Easy]
\begin{enumerate}[label=(\alph*),topsep=2pt,itemsep=1pt,partopsep=0pt,parsep=0pt]
    \item {Amplify only DC signals}
    \item {Pass all frequencies equally}
    \item {Reject a specific range of frequencies}
    \item {Pass only a narrow band of frequencies}
\end{enumerate}
\vspace{2mm}

\noindent\textbf{Q19. Compared to passive filters, active filters generally offer:} \hfill [Topic: active filter | Easy]
\begin{enumerate}[label=(\alph*),topsep=2pt,itemsep=1pt,partopsep=0pt,parsep=0pt]
    \item {Lower cost always}
    \item {No need for a power supply}
    \item {Gain greater than unity}
    \item {The ability to work only at DC}
\end{enumerate}
\vspace{2mm}

\noindent\textbf{Q20. When used as a comparator, an op-amp output switches based on:} \hfill [Topic: op-amp as a comparator and its application in anti-lock braking systems | Easy]
\begin{enumerate}[label=(\alph*),topsep=2pt,itemsep=1pt,partopsep=0pt,parsep=0pt]
    \item {Which input voltage is larger}
    \item {The capacitor charge time}
    \item {The frequency of the input}
    \item {The resistor tolerance in the circuit, which determines the exact threshold voltage for switching the output state}
\end{enumerate}
\vspace{2mm}

\noindent\textbf{Q21. A passive RC low-pass filter has R = 1 k $\Omega$ and C = 0.1 $\mu$ F . What is its cutoff frequency?} \hfill [Topic: low-pass filter | Medium]
\begin{enumerate}[label=(\alph*),topsep=2pt,itemsep=1pt,partopsep=0pt,parsep=0pt]
    \item {1.00 kHz}
    \item {15.9 kHz}
    \item {159 Hz}
    \item {1.59 kHz}
\end{enumerate}
\vspace{2mm}

\noindent\textbf{Q22. In a passive RC high-pass filter, how does the output behave as the input frequency approaches zero (DC)?} \hfill [Topic: high-pass filter | Medium]
\begin{enumerate}[label=(\alph*),topsep=2pt,itemsep=1pt,partopsep=0pt,parsep=0pt]
    \item {Output becomes infinite}
    \item {Output approaches zero}
    \item {Output doubles the input}
    \item {Output equals the input}
\end{enumerate}
\vspace{2mm}

\noindent\textbf{Q23. A band-pass filter has a lower cutoff of 200 Hz and an upper cutoff of 2000 Hz . What is its bandwidth?} \hfill [Topic: band-pass filter | Medium]
\begin{enumerate}[label=(\alph*),topsep=2pt,itemsep=1pt,partopsep=0pt,parsep=0pt]
    \item {632 Hz}
    \item {2200 Hz}
    \item {1000 Hz}
    \item {1800 Hz}
\end{enumerate}
\vspace{2mm}

\noindent\textbf{Q24. An inverting amplifier uses R f = 100 k $\Omega$ and R in = 10 k $\Omega$ . What is its voltage gain?} \hfill [Topic: inverting and non-inverting op-amp | Medium]
\begin{enumerate}[label=(\alph*),topsep=2pt,itemsep=1pt,partopsep=0pt,parsep=0pt]
    \item {- 11}
    \item {+ 10}
    \item {- 10}
    \item {+ 0.1}
\end{enumerate}
\vspace{2mm}

\noindent\textbf{Q25. A non-inverting amplifier has R f = 90 k $\Omega$ and R 1 = 10 k $\Omega$ . What is its voltage gain?} \hfill [Topic: inverting and non-inverting op-amp | Medium]
\begin{enumerate}[label=(\alph*),topsep=2pt,itemsep=1pt,partopsep=0pt,parsep=0pt]
    \item {+ 9}
    \item {- 9}
    \item {+ 10}
    \item {- 10}
\end{enumerate}
\vspace{2mm}

\noindent\textbf{Q26. In an inverting op-amp circuit, why is the inverting input terminal called a virtual ground?} \hfill [Topic: simple op amp circuits --- virtual ground concept | Medium]
\begin{enumerate}[label=(\alph*),topsep=2pt,itemsep=1pt,partopsep=0pt,parsep=0pt]
    \item {It is directly wired to the ground rail}
    \item {It carries large current to ground}
    \item {It floats at the supply voltage}
    \item {It is held at nearly 0 V by feedback without a physical ground connection}
\end{enumerate}
\vspace{2mm}

\noindent\textbf{Q27. For an ideal operational amplifier, what are the input impedance and output impedance respectively?} \hfill [Topic: properties of operational amplifier | Medium]
\begin{enumerate}[label=(\alph*),topsep=2pt,itemsep=1pt,partopsep=0pt,parsep=0pt]
    \item {Infinite and infinite}
    \item {Infinite and zero}
    \item {Zero and zero}
    \item {Zero and infinite}
\end{enumerate}
\vspace{2mm}

\noindent\textbf{Q28. An op-amp has an open-loop gain of 200 , 000 and a maximum output of $\pm$ 12 V . What differential input voltage drives the output to saturation?} \hfill [Topic: device properties of the operational amplifier | Medium]
\begin{enumerate}[label=(\alph*),topsep=2pt,itemsep=1pt,partopsep=0pt,parsep=0pt]
    \item {60 $\mu$ V}
    \item {60 mV}
    \item {600 $\mu$ V}
    \item {6 mV}
\end{enumerate}
\vspace{2mm}

\noindent\textbf{Q29. A summing amplifier has inputs 1 V and 2 V through equal input resistors R , with feedback resistor also R . What is the output voltage?} \hfill [Topic: op-amp as an adder and subtractor | Medium]
\begin{enumerate}[label=(\alph*),topsep=2pt,itemsep=1pt,partopsep=0pt,parsep=0pt]
    \item {- 6 V}
    \item {+ 3 V}
    \item {- 3 V}
    \item {- 1.5 V}
\end{enumerate}
\vspace{2mm}

\noindent\textbf{Q30. In a basic difference amplifier with all four resistors equal, if V 1 = 3 V (non-inverting path) and V 2 = 1 V (inverting path), what is the output?} \hfill [Topic: op-amp as an adder and subtractor | Medium]
\begin{enumerate}[label=(\alph*),topsep=2pt,itemsep=1pt,partopsep=0pt,parsep=0pt]
    \item {- 2 V}
    \item {- 4 V}
    \item {+ 2 V}
    \item {+ 4 V}
\end{enumerate}
\vspace{2mm}

\noindent\textbf{Q31. An op-amp integrator is given a constant positive DC input voltage. What is the resulting output waveform?} \hfill [Topic: op-amp RC circuits --- op-amp integrator | Medium]
\begin{enumerate}[label=(\alph*),topsep=2pt,itemsep=1pt,partopsep=0pt,parsep=0pt]
    \item {A rising exponential}
    \item {A square wave}
    \item {A linearly decreasing (negative-going) ramp}
    \item {A constant positive level}
\end{enumerate}
\vspace{2mm}

\noindent\textbf{Q32. An op-amp differentiator receives a linear ramp input. What is the output?} \hfill [Topic: op-amp differentiator | Medium]
\begin{enumerate}[label=(\alph*),topsep=2pt,itemsep=1pt,partopsep=0pt,parsep=0pt]
    \item {A linear ramp}
    \item {A constant DC level}
    \item {An exponential decay}
    \item {A sine wave}
\end{enumerate}
\vspace{2mm}

\noindent\textbf{Q33. An active high-pass filter uses R = 10 k $\Omega$ and C = 10 nF in its frequency-selective network. What is the cutoff frequency?} \hfill [Topic: active high pass filter | Medium]
\begin{enumerate}[label=(\alph*),topsep=2pt,itemsep=1pt,partopsep=0pt,parsep=0pt]
    \item {15.9 kHz}
    \item {3.18 kHz}
    \item {1.59 kHz}
    \item {159 Hz}
\end{enumerate}
\vspace{2mm}

\noindent\textbf{Q34. A first-order active band-pass filter is commonly formed by cascading which two stages?} \hfill [Topic: active band pass filter | Medium]
\begin{enumerate}[label=(\alph*),topsep=2pt,itemsep=1pt,partopsep=0pt,parsep=0pt]
    \item {A comparator and an integrator}
    \item {Two identical high-pass stages}
    \item {Two identical low-pass stages}
    \item {A low-pass stage and a high-pass stage}
\end{enumerate}
\vspace{2mm}

\noindent\textbf{Q35. What is the primary function of an active band-stop (notch) filter?} \hfill [Topic: active band stop filter | Medium]
\begin{enumerate}[label=(\alph*),topsep=2pt,itemsep=1pt,partopsep=0pt,parsep=0pt]
    \item {Pass only a narrow band of frequencies}
    \item {Pass all frequencies above a cutoff}
    \item {Attenuate a specific band while passing frequencies above and below it}
    \item {Attenuate all frequencies equally}
\end{enumerate}
\vspace{2mm}

\noindent\textbf{Q36. What is a key advantage of an active filter over an equivalent passive RC filter?} \hfill [Topic: active filter | Medium]
\begin{enumerate}[label=(\alph*),topsep=2pt,itemsep=1pt,partopsep=0pt,parsep=0pt]
    \item {It only uses inductors and capacitors}
    \item {It can provide gain and avoids bulky inductors}
    \item {It consumes no power}
    \item {It works without any power supply}
\end{enumerate}
\vspace{2mm}

\noindent\textbf{Q37. When an op-amp is used as a comparator with V in slightly greater than the reference voltage V re f , what is the output?} \hfill [Topic: op-amp as a comparator and its application in anti-lock braking systems | Medium]
\begin{enumerate}[label=(\alph*),topsep=2pt,itemsep=1pt,partopsep=0pt,parsep=0pt]
    \item {Negative saturation voltage}
    \item {Zero volts}
    \item {Positive saturation voltage}
    \item {Exactly V re f}
\end{enumerate}
\vspace{2mm}

\noindent\textbf{Q38. In an anti-lock braking system (ABS), how is an op-amp comparator typically used?} \hfill [Topic: op-amp as a comparator and its application in anti-lock braking systems | Medium]
\begin{enumerate}[label=(\alph*),topsep=2pt,itemsep=1pt,partopsep=0pt,parsep=0pt]
    \item {To amplify the brake pedal force linearly}
    \item {To filter out engine noise from the audio system}
    \item {To integrate the vehicle's acceleration over time}
    \item {To compare a wheel's speed signal against a reference and detect impending lock-up}
\end{enumerate}
\vspace{2mm}

\noindent\textbf{Q39. In a practical op-amp integrator, a large feedback resistor R f is placed across the capacitor mainly to:} \hfill [Topic: op-amp RC circuits --- op-amp integrator | Medium]
\begin{enumerate}[label=(\alph*),topsep=2pt,itemsep=1pt,partopsep=0pt,parsep=0pt]
    \item {Convert it into a differentiator}
    \item {Prevent output drift and saturation from DC offset}
    \item {Reduce the input impedance to zero}
    \item {Increase the integration speed}
\end{enumerate}
\vspace{2mm}

\noindent\textbf{Q40. An op-amp has a slew rate of 1 V / $\mu$ s . What is the minimum time required for its output to swing from 0 V to 5 V ?} \hfill [Topic: device properties of the operational amplifier | Medium]
\begin{enumerate}[label=(\alph*),topsep=2pt,itemsep=1pt,partopsep=0pt,parsep=0pt]
    \item {1 $\mu$ s}
    \item {10 $\mu$ s}
    \item {5 $\mu$ s}
    \item {0.2 $\mu$ s}
\end{enumerate}
\vspace{2mm}

\noindent\textbf{Q41. A passive first-order RC low-pass filter has R = 10 k $\Omega$ and C = 15.9 nF . At what frequency does the output power drop to half of the maximum, and what is the phase shift there?} \hfill [Topic: low-pass filter | Hard]
\begin{enumerate}[label=(\alph*),topsep=2pt,itemsep=1pt,partopsep=0pt,parsep=0pt]
    \item {f $\approx$ 1 kHz , phase = - 9 0 $^\circ$}
    \item {f $\approx$ 10 kHz , phase = - 4 5 $^\circ$}
    \item {f $\approx$ 100 Hz , phase = - 4 5 $^\circ$}
    \item {f $\approx$ 1 kHz , phase = - 4 5 $^\circ$}
\end{enumerate}
\vspace{2mm}

\noindent\textbf{Q42. A first-order passive RC high-pass filter is cascaded (unbuffered) with an identical stage. Compared to a single stage, how does the overall -3 dB cutoff frequency behave?} \hfill [Topic: high-pass filter | Hard]
\begin{enumerate}[label=(\alph*),topsep=2pt,itemsep=1pt,partopsep=0pt,parsep=0pt]
    \item {It stays exactly at the single-stage f c}
    \item {It shifts higher than the single-stage f c due to loading interaction}
    \item {It becomes independent of R and C}
    \item {It shifts lower than the single-stage f c}
\end{enumerate}
\vspace{2mm}

\noindent\textbf{Q43. A band-pass filter has a center frequency f 0 = 2 kHz and a quality factor Q = 10 . What is its bandwidth and the approximate lower cutoff frequency?} \hfill [Topic: band-pass filter | Hard]
\begin{enumerate}[label=(\alph*),topsep=2pt,itemsep=1pt,partopsep=0pt,parsep=0pt]
    \item {B W = 100 Hz , f L $\approx$ 1.95 kHz}
    \item {B W = 200 Hz , f L $\approx$ 1 kHz}
    \item {B W = 200 Hz , f L $\approx$ 1.9 kHz}
    \item {B W = 20 kHz , f L $\approx$ 1.9 kHz}
\end{enumerate}
\vspace{2mm}

\noindent\textbf{Q44. In an inverting amplifier with an ideal op-amp, the inverting input is called a virtual ground. If a real op-amp has finite open-loop gain A = 1 0 5 and the closed-loop gain magnitude is 100, what is the approximate voltage at the inverting node when V o u t = 5 V ?} \hfill [Topic: simple op amp circuits -- virtual ground concept | Hard]
\begin{enumerate}[label=(\alph*),topsep=2pt,itemsep=1pt,partopsep=0pt,parsep=0pt]
    \item {500 $\mu$ V}
    \item {0 V exactly}
    \item {5 mV}
    \item {50 $\mu$ V}
\end{enumerate}
\vspace{2mm}

\noindent\textbf{Q45. An inverting amplifier uses R in = 2 k $\Omega$ and R f = 20 k $\Omega$ . A non-inverting amplifier is built to give the same gain magnitude. What feedback and input resistor ratio must the non-inverting stage use?} \hfill [Topic: inverting and non-inverting op-amp | Hard]
\begin{enumerate}[label=(\alph*),topsep=2pt,itemsep=1pt,partopsep=0pt,parsep=0pt]
    \item {R f / R 1 = 0.1}
    \item {R f / R 1 = 11}
    \item {R f / R 1 = 10}
    \item {R f / R 1 = 9}
\end{enumerate}
\vspace{2mm}

\noindent\textbf{Q46. A summing amplifier has inputs V 1 = 1 V through 10 k $\Omega$ and V 2 = 2 V through 20 k $\Omega$ , with R f = 40 k $\Omega$ . What is V o u t ?} \hfill [Topic: op-amp as an adder and subtractor | Hard]
\begin{enumerate}[label=(\alph*),topsep=2pt,itemsep=1pt,partopsep=0pt,parsep=0pt]
    \item {- 12 V}
    \item {- 6 V}
    \item {- 8 V}
    \item {- 4 V}
\end{enumerate}
\vspace{2mm}

\noindent\textbf{Q47. An ideal op-amp integrator has R = 100 k $\Omega$ and C = 1 $\mu$ F . A constant input of 2 V is applied starting at t = 0 with zero initial capacitor charge. What is V o u t at t = 0.5 s ?} \hfill [Topic: op-amp RC circuits -- op-amp integrator | Hard]
\begin{enumerate}[label=(\alph*),topsep=2pt,itemsep=1pt,partopsep=0pt,parsep=0pt]
    \item {- 5 V}
    \item {- 10 V}
    \item {- 20 V}
    \item {- 2 V}
\end{enumerate}
\vspace{2mm}

\noindent\textbf{Q48. A practical op-amp differentiator adds a series input resistor R 1 to the capacitor to limit high-frequency gain. If C = 10 nF , R 1 = 1 k $\Omega$ and R f = 10 k $\Omega$ , above what frequency does the circuit stop differentiating?} \hfill [Topic: op-amp differentiator | Hard]
\begin{enumerate}[label=(\alph*),topsep=2pt,itemsep=1pt,partopsep=0pt,parsep=0pt]
    \item {$\approx$ 15.9 kHz}
    \item {$\approx$ 100 kHz}
    \item {$\approx$ 1.59 kHz}
    \item {$\approx$ 159 kHz}
\end{enumerate}
\vspace{2mm}

\noindent\textbf{Q49. A first-order active high-pass filter (non-inverting) uses R = 15.9 k $\Omega$ , C = 10 nF , and sets a passband gain of 11. What are its cutoff frequency and passband gain in dB?} \hfill [Topic: active high pass filter | Hard]
\begin{enumerate}[label=(\alph*),topsep=2pt,itemsep=1pt,partopsep=0pt,parsep=0pt]
    \item {f c $\approx$ 1 kHz , $\approx$ 20.8 dB}
    \item {f c $\approx$ 100 Hz , $\approx$ 21 dB}
    \item {f c $\approx$ 10 kHz , $\approx$ 20.8 dB}
    \item {f c $\approx$ 1 kHz , $\approx$ 11 dB}
\end{enumerate}
\vspace{2mm}

\noindent\textbf{Q50. An active band-pass filter is formed by cascading a low-pass stage with f cH = 8 kHz and a high-pass stage with f c L = 500 Hz . What is the approximate center frequency (geometric mean) and bandwidth?} \hfill [Topic: active band pass filter | Hard]
\begin{enumerate}[label=(\alph*),topsep=2pt,itemsep=1pt,partopsep=0pt,parsep=0pt]
    \item {f 0 = 4.25 kHz , B W = 7.5 kHz}
    \item {f 0 = 2 kHz , B W = 8.5 kHz}
    \item {f 0 = 2 kHz , B W = 7.5 kHz}
    \item {f 0 = 4 kHz , B W = 500 Hz}
\end{enumerate}
\vspace{2mm}

\noindent\textbf{Q51. A band-stop (notch) filter is required to reject 50 Hz mains interference with a quality factor Q = 25 . What is the rejection bandwidth?} \hfill [Topic: active band stop filter | Hard]
\begin{enumerate}[label=(\alph*),topsep=2pt,itemsep=1pt,partopsep=0pt,parsep=0pt]
    \item {1250 Hz}
    \item {25 Hz}
    \item {0.5 Hz}
    \item {2 Hz}
\end{enumerate}
\vspace{2mm}

\noindent\textbf{Q52. In an ABS system, a comparator compares a wheel-speed voltage against a reference proportional to vehicle speed. Hysteresis is added using positive feedback. What is the primary benefit of this hysteresis?} \hfill [Topic: op-amp as a comparator and its application in anti-lock braking systems | Hard]
\begin{enumerate}[label=(\alph*),topsep=2pt,itemsep=1pt,partopsep=0pt,parsep=0pt]
    \item {It converts the comparator into a linear amplifier}
    \item {It eliminates the need for a reference voltage}
    \item {It prevents rapid chattering of the output near the threshold due to noise}
    \item {It increases the comparator's open-loop gain}
\end{enumerate}
\vspace{2mm}

\noindent\textbf{Q53. An op-amp has a gain-bandwidth product (GBW) of 10 MHz . It is configured as a non-inverting amplifier with a gain of 50. What is the approximate closed-loop bandwidth?} \hfill [Topic: properties of operational amplifier | Hard]
\begin{enumerate}[label=(\alph*),topsep=2pt,itemsep=1pt,partopsep=0pt,parsep=0pt]
    \item {50 kHz}
    \item {500 kHz}
    \item {10 MHz}
    \item {200 kHz}
\end{enumerate}
\vspace{2mm}

\noindent\textbf{Q54. An op-amp has a slew rate of 0.5 V / $\mu$ s . What is the maximum frequency of a sine wave with 10 V peak amplitude that can be reproduced without slew-rate distortion?} \hfill [Topic: operational amplifier abstraction- device properties of the operational amplifier | Hard]
\begin{enumerate}[label=(\alph*),topsep=2pt,itemsep=1pt,partopsep=0pt,parsep=0pt]
    \item {$\approx$ 50 kHz}
    \item {$\approx$ 79.6 kHz}
    \item {$\approx$ 15.9 kHz}
    \item {$\approx$ 7.96 kHz}
\end{enumerate}
\vspace{2mm}

\noindent\textbf{Q55. A difference amplifier has all four resistors equal to R . If V 1 = 3 V is applied to the non-inverting path and V 2 = 5 V to the inverting path, what is V o u t , and what is the ideal common-mode gain?} \hfill [Topic: op-amp as an adder and subtractor | Hard]
\begin{enumerate}[label=(\alph*),topsep=2pt,itemsep=1pt,partopsep=0pt,parsep=0pt]
    \item {- 2 V , common-mode gain = 0}
    \item {- 2 V , common-mode gain = 1}
    \item {+ 8 V , common-mode gain = 0}
    \item {+ 2 V , common-mode gain = 0}
\end{enumerate}
\vspace{2mm}

\noindent\textbf{Q56. A practical op-amp integrator places a resistor R f = 1 M $\Omega$ in parallel with the feedback capacitor C = 0.1 $\mu$ F , with input resistor R = 10 k $\Omega$ . What is the DC gain and why is R f needed?} \hfill [Topic: op-amp RC circuits -- op-amp integrator | Hard]
\begin{enumerate}[label=(\alph*),topsep=2pt,itemsep=1pt,partopsep=0pt,parsep=0pt]
    \item {DC gain = - 100 ; R f limits DC gain to prevent output saturation from offset/bias}
    \item {DC gain = - 1000 ; R f sets the unity-gain frequency}
    \item {DC gain = 0 ; R f provides AC coupling only}
    \item {DC gain = - 10 ; R f increases integration accuracy at high frequency}
\end{enumerate}
\vspace{2mm}

\noindent\textbf{Q57. A Sallen-Key second-order low-pass filter is designed for a Butterworth response. What damping/quality factor and cutoff behavior characterize this maximally flat response?} \hfill [Topic: active filter | Hard]
\begin{enumerate}[label=(\alph*),topsep=2pt,itemsep=1pt,partopsep=0pt,parsep=0pt]
    \item {Q = 0.5 , critically damped with 20 dB/decade rolloff}
    \item {Q = 1 , with a peaking passband and 20 dB/decade rolloff}
    \item {Q = 1/ 2 $\approx$ 0.707 , maximally flat passband with 40 dB/decade rolloff}
    \item {Q = 2 , with ripple in the passband and 40 dB/decade rolloff}
\end{enumerate}
\vspace{2mm}

\noindent\textbf{Q58. An ideal op-amp differentiator with R f = 10 k $\Omega$ and C = 1 $\mu$ F receives a triangular input that ramps linearly from 0 to 4 V over 2 ms . What is the output voltage magnitude during this ramp?} \hfill [Topic: op-amp differentiator | Hard]
\begin{enumerate}[label=(\alph*),topsep=2pt,itemsep=1pt,partopsep=0pt,parsep=0pt]
    \item {40 mV}
    \item {200 mV}
    \item {20 mV}
    \item {2 V}
\end{enumerate}
\vspace{2mm}

\noindent\textbf{Q59. An inverting amplifier is designed for a gain of -20 with R f = 200 k $\Omega$ . The op-amp has an input bias current. To minimize DC output error from bias currents, what resistor should be placed at the non-inverting input?} \hfill [Topic: inverting and non-inverting op-amp | Hard]
\begin{enumerate}[label=(\alph*),topsep=2pt,itemsep=1pt,partopsep=0pt,parsep=0pt]
    \item {10 k $\Omega$ equal to R in}
    \item {0 $\Omega$ (direct ground connection)}
    \item {200 k $\Omega$ equal to R f}
    \item {$\approx$ 9.5 k $\Omega$ (parallel combination of R in and R f )}
\end{enumerate}
\vspace{2mm}

\noindent\textbf{Q60. A comparator-based ABS controller detects impending wheel lock when wheel deceleration exceeds a threshold. If the wheel-speed sensor outputs a voltage proportional to speed and the circuit differentiates then compares it, what condition triggers the brake-release signal?} \hfill [Topic: op-amp as a comparator and its application in anti-lock braking systems | Hard]
\begin{enumerate}[label=(\alph*),topsep=2pt,itemsep=1pt,partopsep=0pt,parsep=0pt]
    \item {When the raw wheel-speed voltage equals the supply voltage}
    \item {When the comparator operates in its linear region}
    \item {When the differentiated signal (deceleration) crosses the negative reference threshold}
    \item {When the integrated speed signal reaches zero}
\end{enumerate}
\vspace{2mm}


\newpage
\section{Section B: Comprehensive Analytical Problems with Step-by-Step Solutions}

\begin{subjectiveq}{5}{1}
An operational amplifier IC 741 with supply rails $V_{CC} = \pm 15\text{ V}$ is utilized in the following configurations:\begin{enumerate}[label=(\alph*)]\item Design an inverting summing amplifier to produce an output voltage $v_o = -(3v_1 + 2v_2 + 0.5v_3)$ using an input feedback resistor $R_f = 60\text{ k}\Omega$.\item In an active first-order low-pass Butterworth filter with a passband voltage gain of $1.586$ and cutoff frequency $f_c = 2.5\text{ kHz}$, calculate the component values of $R$ and $C$ (choosing $C = 10\text{ nF}$) and feedback resistors $R_1, R_f$.\item Formulate how an op-amp comparator serves as the threshold detector in an automotive Anti-Lock Braking System (ABS) wheel-slip detection unit.\end{enumerate}
\end{subjectiveq}

\begin{solbox}
\textbf{Step 1: Inverting Summing Amplifier Design}\\Output equation: $v_o = -\left( \frac{R_f}{R_1}v_1 + \frac{R_f}{R_2}v_2 + \frac{R_f}{R_3}v_3 \right)$.With $R_f = 60\text{ k}\Omega$:\begin{align*}R_1 &= \frac{60\text{ k}\Omega}{3} = 20\text{ k}\Omega, \quad R_2 = \frac{60\text{ k}\Omega}{2} = 30\text{ k}\Omega, \quad R_3 = \frac{60\text{ k}\Omega}{0.5} = 120\text{ k}\Omega\end{align*}\textbf{Step 2: Active First-Order Low-Pass Filter Design}\\Given $f_c = 2.5\text{ kHz}$ and $C = 10\text{ nF}$:\[ R = \frac{1}{2\pi f_c C} = \frac{1}{2\pi \times 2500 \times 10 \times 10^{-9}} \approx 6.37\text{ k}\Omega \]For non-inverting passband gain $A_v = 1 + \frac{R_f}{R_1} = 1.586$:\[ \frac{R_f}{R_1} = 0.586 \implies \text{Selecting } R_1 = 10\text{ k}\Omega \implies R_f = 5.86\text{ k}\Omega \]\textbf{Step 3: Anti-lock Braking System (ABS) Comparator Architecture}\\In modern automotive ABS, variable reluctance wheel-speed sensors produce a pulsed AC voltage proportional to wheel speed $\omega_w$.\begin{itemize}[leftmargin=1.5em]\item Inverting terminal receives conditioned sensor signal $v_{sensor}$.\item Non-inverting terminal receives reference threshold $V_{ref}$ from the Electronic Control Unit.\item When wheel slip occurs (impending lockup), sensor voltage drops below $V_{ref}$, causing comparator output to jump to $+V_{sat}$ (logic HIGH).\item The ECU pulses solenoid valves to release and modulate brake caliper pressure, preserving vehicle steering control.\end{itemize}
\end{solbox}

\begin{examinertip}
\textbf{Key Takeaway}: Stating virtual ground ($V^+ \approx V^-$) and zero input bias current ($I_{in} \approx 0$) is vital for full marks in op-amp proofs.
\end{examinertip}
\vspace{4mm}
